\clearpage\SourceLeaf{520}
\renewcommand{\BookRunningTitle}{Section Eight.}
\pdfbookmark[0]{Section Eight. A Fuller Consideration of Certain Investigations Concerning Circle Division.}{sectio-octava}
\begin{center}
\textbf{SECTION EIGHT.}\\[1em]
A FULLER CONSIDERATION OF CERTAIN INVESTIGATIONS CONCERNING CIRCLE DIVISION.
\end{center}
\BookRule{.18\textwidth}
\par\medskip\Needspace{4\baselineskip}\pdfbookmark[1]{367.}{sectio-octava-art-367}\centerline{\textbf{367.}}\nobreak\smallskip
What we set out in the latter part of Section Seven, from Article $355$ onward, certainly gives substantial examples of the great fertility of the theory of circle division, and of the remarkable connection linking this subject with various \textit{arithmetic} investigations. There, however, the severe limits of space and time allowed us only to touch lightly on this field, which rewards our efforts with a richer harvest the further we advance into it. We therefore intend to take up here again one or another of the questions begun there and discuss them more fully: readers will certainly see with great wonder that the solution of several arithmetic problems, which anyone would have expected to be worlds apart from this subject, rests on this foundation.
\par\medskip\Needspace{4\baselineskip}\pdfbookmark[1]{368.}{sectio-octava-art-368}\centerline{\textbf{368.}}\nobreak\smallskip
The investigation begun in Article $356$ offers an especially fruitful subject. There we divided the collection of roots of the equation $x^n-1=0$ (excluding unity) into two classes and showed how to determine the sum in each class. These sums were $=-\frac12+\frac12\surd n$ and $-\frac12-\frac12\surd n$ when $n$ has the form $4n+1$, or $=-\frac12+\frac12\surd -n$ and $-\frac12-\frac12\surd -n$ when $n$ has the form $4n+3$. Yet there we had not only \SourceLeaf{521}restricted ourselves to the case where $n$ is a \textit{prime number}; we had also, what was of much greater importance still, left the \textit{sign} of the radical quantity undetermined, or rather neglected to support this determination, indicated in a few words, with a solid proof. We must therefore make good these deficiencies before anything else.
\par\medskip\Needspace{4\baselineskip}\pdfbookmark[1]{369.}{sectio-octava-art-369}\centerline{\textbf{.369.}}\nobreak\smallskip
Now let $n$ be any positive integer, and $R$ a root of the equation $x^n-1\pdot=0$ none of whose powers preceding the power with exponent $n$ equals unity (see Article $359$, II.). As in Section VII, let us denote by $[\lambda]$ the power $R^\lambda$, so that $[0]=1$, $[1]$, $[2]$, $[3]\dotsfour[n-1]$ represent all the roots of the equation $x^n-1=0$. Further, let us denote the sum
\[
{}[0]+[1]+[4]+[9]+\dotsfour+[(n-1)^2]\quad\text{by }\Sigma[\mathfrak{Q}]
\]
and more generally
\[
{}[0]+[\lambda]+[4\lambda]+[9\lambda]\dotsfour+[\lambda(n-1)^2]\quad\text{by }\Sigma[\mathfrak{Q}\lambda]
\]
so that $\mathfrak{Q}$ denotes, without further specification, the squares of the numbers $0,1,2,3\dotsfour n-1$. Thus, just as in general $[\lambda]=[\mu]$ when $\lambda$, $\mu$ are any integers (positive or negative) congruent modulo $n$, we also have $\Sigma[\mathfrak{Q}\lambda]=\Sigma[\mathfrak{Q}\mu]$ when $\lambda\equiv\mu$. With these preparations we have the following.
\par\medskip\Needspace{4\baselineskip}\pdfbookmark[1]{370.}{sectio-octava-art-370}\centerline{\textbf{370.}}\nobreak\smallskip
\textsc{Problem.} \textit{Determine the product of the two sums $\Sigma[\mathfrak{Q}]$ and $\Sigma[-\mathfrak{Q}]$.}
\par\noindent\textit{Solution.} Since $nn\equiv0$, $(n+1)^2\equiv1$, $(n+2)^2\equiv4$, etc. $(\mathrm{mod}.n)$, it is easy to see that $\Sigma[\mathfrak{Q}]$ becomes
\[
\begin{aligned}
&=[1]+[4]+[9]+[16]\dotsfive+[nn]\\
&=[4]+[9]+[16]+[25]\dotsfive+[(n+1)]^2\\
&=[9]+[16]+[25]+[36]\dotsfive+[(n+2)^2]
\end{aligned}
\]
\centerline{etc.\quad or in general}
\[
{}=[kk]+[(k+1)^2]+[(k+2)^2]+[(k+3)^2]\dotsfour+[(n+k-1)^2]
\]
Hence $[-kk]\times\Sigma[\mathfrak{Q}]$
\[
{}=[0]+[2k+1]+[4k+4]+[6k+9]\dotsfive+[(n-1)^2+2(n-1)k]
\]
